% =====================================================================
%  macros.tex — comandos del informe GA-EXP-02 (no hace falta editarlo)
% =====================================================================

% Texto pendiente: aparece en ROJO en el PDF para que no se olvide nada.
\newcommand{\completar}[1]{\textcolor{red}{\textbf{[COMPLETAR: #1]}}}
\newcommand{\faltanivel}{\textcolor{red}{\textbf{?}}}

% ---------- Identificación del equipo evaluado ----------
% \identificacion{Equipo}{Tema}{Fecha}{Hora inicio}{Hora fin}
\newcommand{\identificacion}[5]{%
  \subsection*{Identificación}
  \begin{tabular}{@{}p{0.28\textwidth}p{0.66\textwidth}@{}}
    \toprule
    Equipo evaluado & #1 \\
    Tema            & #2 \\
    Fecha           & #3 \\
    Hora de inicio  & #4 \\
    Hora de fin     & #5 \\
    \bottomrule
  \end{tabular}
}

% ---------- Tabla de dimensiones ----------
\newenvironment{dimensiones}{%
  \subsection*{Evaluación por dimensión}
  \renewcommand{\arraystretch}{1.25}
  \begin{longtable}{@{}>{\raggedright\arraybackslash}p{0.24\textwidth}>{\centering\arraybackslash}p{0.07\textwidth}>{\raggedright\arraybackslash}p{0.63\textwidth}@{}}
    \toprule
    \textbf{Dimensión (peso)} & \textbf{Nivel} & \textbf{Evidencia observable} \\
    \midrule
    \endhead
}{%
    \bottomrule
  \end{longtable}
}
% \dimension{Nombre (peso)}{Nivel 0-4}{Evidencia}
\newcommand{\dimension}[3]{#1 & #2 & #3 \\ \midrule}

% ---------- Nota sugerida con la cuenta escrita ----------
% \notasugerida{D1}{D2}{D3}{D4}{D5}{D6}   (niveles 0-4)
% Fórmula de la guía: (10*D1 + 20*D2 + 15*D3 + 20*D4 + 25*D5 + 10*D6) / 4 * 0.1
\newcommand{\notasugerida}[6]{%
  \subsection*{Nota sugerida}
  \edef\NSsuma{\fpeval{10*(#1)+20*(#2)+15*(#3)+20*(#4)+25*(#5)+10*(#6)}}%
  \edef\NSdiv{\fpeval{\NSsuma/4}}%
  \edef\NSnota{\fpeval{\NSdiv*0.1}}%
  Fórmula: $\dfrac{10D_1+20D_2+15D_3+20D_4+25D_5+10D_6}{4}\times 0.1$
  \[
    \frac{10(#1)+20(#2)+15(#3)+20(#4)+25(#5)+10(#6)}{4}\times 0.1
    = \frac{\NSsuma}{4}\times 0.1 = \NSdiv\times 0.1 = \NSnota
    \approx \mathbf{\fpeval{round(\NSnota,1)}}
  \]
}

% ---------- Retroalimentación ----------
% \retroalimentacion{Fortaleza}{Mejora}
\newcommand{\retroalimentacion}[2]{%
  \subsection*{Fortaleza}
  #1
  \subsection*{Mejora accionable}
  #2
}
